Une solution de chlorure de potassium KCl a une concentration CKCl= 5 .10-3 mol.L-1 .

a)Ecrire la réaction de dissolution dans l'eau du chlorure de potassium .

b)La dissolution est totale . Calculer , en mol.m-3 , les concentrations dans la solution des ions K+ et Cl- ? Justifiez clairement votre réponse .

c)On donne les conductivités molaires ioniques : l(Cl-)=7,6 10-3 S.m2.mol-1 et l(K+)= 7,4 10-3 S.m2.mol-1 . Calculer la conductivité de la solution .
